\chapter{Fase 2: limpiar y ordenar los datos (Silver y Gold)}
\label{cap:d-f2}
\textbf{Fechas}: del 28 de septiembre al 2 de octubre de 2026. \textbf{Notas}: \codigo{04-silver.md},
\codigo{10-silver-fase5.md}, \codigo{18-cierre-fase-2.md}.

\textbf{Qué se buscaba}: convertir 14 fuentes distintas en tablas limpias, con las mismas reglas, documentadas y
consultables.

\textbf{Por qué duró tanto}: Brandon decidió construir la página en paralelo y abrir las Fases 3 y 4 con las fuentes que ya
estaban limpias. Las demás se limpiaron justo antes de la fase que las usaba: clima, huracanes y tipo de cambio antes del
Pronóstico; nacionalidades, vuelos, cruceros y reseñas antes de la campaña. Una auditoría (29-sep) lo hizo visible y la fase
se cerró el 2 de octubre.

\begin{opciones}[2.1 · La página, en paralelo · 28-sep-2026]
\textbf{Elegida}: construir la página web desde la Fase 2, junto con los datos. \textbf{Descartado}: esperar al final.
\textbf{Consecuencia}: la página creció con cada fase (ver el capítulo~\ref{cap:d-pagina}).
\end{opciones}

\begin{opciones}[2.2 · Qué es ``oferta turística'' · 28-sep-2026]
\textbf{Elegida}: los giros SCIAN característicos del turismo según SECTUR e INEGI: alojamiento (721), alimentos y bebidas
(722), agencias de viajes (5615), transporte turístico (487), museos y sitios (712) y esparcimiento (713).

\textbf{Resultado}: 875,667 negocios turísticos en el país (14.3\,\%) y \textbf{13,663 en Quintana Roo} (19.7\,\% de sus
69,260).
\end{opciones}

\begin{opciones}[2.3 · Tren Maya con dos renglones por mes · 28-sep-2026]
\textbf{Elegida}: sumar las estaciones de cada destino. \textbf{Evidencia}: la suma cuadra exacto con el total de la zona.
Enero de 2025, Gran Costa Maya: $7{,}084$ = Chetumal ($24+3{,}478$) + Bacalar ($213+3{,}369$).
\end{opciones}

\begin{opciones}[2.4 · Los ceros imposibles · 28-sep-2026]
\textbf{La pregunta}: SITUR-Q trae ceros donde no puede haber cero.

\begin{center}
\small
\begin{tabularx}{\linewidth}{>{\raggedright\arraybackslash}p{0.3\linewidth}YY}
\encabezado \h{Caso} & \h{Regla} & \h{Por qué} \\
Ocupación con 0 cuartos disponibles & Hueco (693 filas) & Físicamente imposible \\
\filagris Afluencia y derrama en 0 & Hueco & Un destino no recibe cero turistas ni cero pesos \\
Los 4 aeropuertos en 0 todo 2025 & Hueco si \emph{todos} marcan 0 ese mes (114 filas) & Cancún promedió 72.2\,\% de
ocupación en 2025: es la fuente la que dejó de publicar \\
\filagris Zona arqueológica con 0 & Se conserva, marcada & Suele ser un cierre real (Muyil, 2024–2026) \\
\filagris Cruceristas de 2020 en 0 & Se conserva & Puertos cerrados por la pandemia \\
\bottomrule
\end{tabularx}
\normalsize
\end{center}
\textbf{Descartado}: conservar los ceros. La página habría dicho ``0 pasajeros en avión en 2025''. \textbf{Consecuencia}: la
página usa 2024 como último año completo del avión (15,959,277 pasajeros) y lo dice.
\end{opciones}

\begin{opciones}[2.5 · Qué tormenta ``afecta al sur'' · 30-sep-2026]
\begin{center}
\small
\begin{tabularx}{\linewidth}{>{\raggedright\arraybackslash}p{0.34\linewidth}rrY}
\encabezado \h{Opción} & \h{Eventos} & \h{Por año} & \h{Problema} \\
\textbf{200 km · ≥ 34 nudos · desde 1966 (elegida)} & \textbf{31 en 60 años} & \textbf{0.517} & — \\
100 km · ≥ 34 nudos · desde 1966 & 13 & 0.217 & Muy pocos para una tasa por mes \\
300 km · ≥ 34 nudos · desde 1966 & 59 & 0.983 & Mete tormentas sin viento fuerte en el sur \\
200 km · ≥ 34 nudos · desde 1851 & 82 en 175 & 0.469 & Antes de los satélites se perdían tormentas \\
\end{tabularx}
\normalsize
\end{center}
\textbf{Evidencia}: 23 de los 31 eventos caen de agosto a octubre; entre ellos Carmen (1974, a 15.3 km de Chetumal con 125
nudos) y Dean (2007). Se guardaron los 55,524 puntos de trayectoria con su distancia, para poder cambiar el radio sin
volver al texto crudo.
\end{opciones}

\begin{opciones}[2.6 · Clima: dos tablas y hora local · 30-sep-2026]
\textbf{Decisión técnica}: una tabla diaria (1950–2026, para medir la temporada de lluvias) y una horaria (2019–2026, para
la Torre). El crudo horario viene en UTC; se pasa a hora local restando 5 horas (Quintana Roo usa UTC$-5$ fijo desde 2015).
0 huecos en los 128 archivos.

\textbf{Hueco declarado}: no hay punto de clima propio para la Laguna Milagros ni para Calderitas; se usa Chetumal (a 7.7 km
de Calderitas y 18.9 km de Xul-Ha), y se dice.
\end{opciones}

\begin{opciones}[2.7 · Feriados sin tipo de cambio · 30-sep-2026]
\textbf{Elegida}: \textbf{se quedan vacíos} (336 días), con su bandera. El promedio del mes se saca solo con los días
observados y se guarda cuántos fueron (agosto de 2026: 21 días, 17.0609 pesos por dólar).

\textbf{Descartado}: copiar el día anterior, porque inventa una cotización que no existió; y usar el último dato del mes,
porque depende de un solo día.
\end{opciones}

\begin{opciones}[2.8 · AFAC: dos aerolíneas con la misma etiqueta · 2-oct-2026]
\textbf{La pregunta}: el archivo de vuelos traía 300 llaves repetidas: 98 renglones idénticos (se deja uno), 177 con una
cifra y una copia en cero (se conserva la cifra) y \textbf{25 con dos cifras distintas}, todas de ``Virgin America (Alaska
Airlines)'', de 2016 a enero de 2018.

\begin{center}
\small
\begin{tabularx}{\linewidth}{>{\raggedright\arraybackslash}p{0.3\linewidth}Y}
\encabezado \h{Opción} & \h{Por qué sí o por qué no} \\
\textbf{Sumarlas (elegida)} & Son dos aerolíneas que se fusionaron en 2018 bajo una sola etiqueta; en juego hay 369,311
pasajeros de 1,042 millones (0.035\,\%) \\
\filagris Dos renglones marcados & Obliga a recordarlo en cada suma \\
Dejar AFAC fuera & Se pierde una fuente por 25 renglones \\
\bottomrule
\end{tabularx}
\normalsize
\end{center}
\textbf{Regla de paro}: si aparece otra etiqueta con dos cifras, el proceso se detiene.
\end{opciones}

\begin{opciones}[2.9 · Dónde consultar los datos: el almacén · 2-oct-2026]
\textbf{Elegida}: \textbf{DuckDB}, una base en un solo archivo que lee los Parquet sin copiarlos, con un \textbf{modelo
estrella}: la tabla de hechos (lugar $\times$ mes $\times$ variable, 3,287 renglones) y las dimensiones de lugares (15) y
meses (192). \textbf{Un hueco no tiene renglón}.

\textbf{Descartado}: PostgreSQL, porque exige un servidor y el proyecto debe correr sin internet en cualquier computadora.

\textbf{Consecuencia}: un diccionario de datos generado solo, de \textbf{56 tablas}. Si una columna no tiene significado
escrito, el programa se detiene.
\end{opciones}

\begin{opciones}[2.10 · Cuando dos fuentes no coinciden]
\textbf{Isla Mujeres}: las noticias decían 93.5\,\% y 95\,\% de ocupación en 2026; DataTur registra 74.7\,\% (febrero) y
43.2\,\% (abril). \textbf{Regla}: el dato oficial va antes que la nota de prensa.

\textbf{Cruceros}: en Mahahual, SITUR-Q cuenta siempre entre 11\,\% y 35\,\% más que DataTur; en Cozumel casi coinciden
(2025: $4{,}915{,}242\div4{,}724{,}255-1=+4.0\,\%$). \textbf{Regla}: una fuente por lugar, y la diferencia se declara.
\end{opciones}

\section*{Hallazgo para la campaña}
Al aeropuerto de Chetumal llegaron \textbf{250 extranjeros} en todo 2025; a Cancún, 9,408,423. El viajero extranjero no llega
al sur en avión: el anuncio tiene que alcanzarlo antes del viaje o cuando ya está en Cancún. Esto decidió las personas de la
Fase 8.

\begin{error}
\begin{itemize}
\item \textbf{Afluencia y derrama} parecían llegar a junio de 2024 por unos ceros; terminan en marzo de 2024.
\item \textbf{Huracanes}: dos líneas mal escritas en el archivo oficial. Una (1975) no dice si la latitud es norte o sur: se
  dejó vacía, sin adivinar. La primera exploración saltaba esa línea sin avisar (55,523 puntos); el código final se detiene
  si no entiende una línea (55,524).
\item \textbf{INAH}: el bloque de extranjeros de septiembre de 2025 venía dos veces, una con cifras y otra en ceros.
\item \textbf{DataTur}: un mismo centro contaba dos veces por las marcas de nota (``ISLA MUJERES /4'') y las notas al pie se
  leían como centros.
\item \textbf{Windows}: la ruta de la carpeta tiene espacios y Spark la cortaba; se usa la ruta corta.
\item \textbf{Revisión del 2 de octubre}: el diccionario se detuvo porque 12 tablas nuevas no tenían descripción.
\end{itemize}
\end{error}
